% Nouveau EGA V (cote 25): archive cote 25, site page 67 (PDF page 68).
% typescript. Draft compared: none.
% AI-transcribed by Claude Opus 5.5 (claude-opus-5-5, adaptive thinking, effort high), Message Batches API, batch ega25b (msgbatch_01A7n5YuXubUPNK1zsepi8RG), 2026-11-13; not yet checked by eye.
\scanpage{25/67}
\mspage{\unsure{56}}

Le préschémad'incidence  (resp. son image inverse dans T) n'est alors
autre que les préschémas des zéros de l'un ou l'autre homomorphisme,
ou encore de lasection $\phi^{m}$ .  On pourrait dénoter le préschéma
d'incidence par $H^{(m)}$ , pour m=1 on retouve celui du n°1 .  Si X est
au dessus de P, on pourra poser $\underline{Y}^{(m)} = X\times_{P} \underline{H}^{(m)}$ ,  et définir de même
la notation $\underline{Y}^{m}$ lorsque \hand{$\zeta$} est un point de $\unsure{\mathrm{Gr}}^{m}$ à valeurs dans un S'.
Ainsi les \xout{\unsure{Y}} $Y^{m}_{\hand{$\zeta$}}$  sont les``\uline{sections} \uline{linéaires}''  de X \inserted{sur P} (ou plutôt des
$X_{S'}$\inserted{sur $P_{S'}$})   par les sous-variétés linéaires de codimension m de P (ou
plutôt des$P_{S'}$).

\margin{$|$}Je profite de l'occasion pour une\xout{xx}  autocritique notationnelle,
qui pourrait venir dès le n°1. C'est le fait  de faire correspondre
arbitrairement, pour indiquer un objet $\underline{Y}$ qui correspond à X, la lettre
Y à X  (de sorte que si X devient Z, on ne sait plus très bien quoi\xout{x}
prendre). Cet inconvénient m'a déjà amené à quelques incohérences de
notations passagères.  Peut-être que le contexte un peu plus général
avec un entier m comme ici suggère une solution raisonnnable: écrire
$\underline{X}^{(m)}$ au lieu de  $\underline{Y}^{(m)}$ , donc  $\underline{X}^{(1)}$ au lieu du  $\underline{Y}$  du n° 1. De cette
façon, on  aurait bien , à peu de choses près ... ,  $\underline{X}^{(m)(m')} = X^{(m+m')}$\hand{.}
Je vais essayer cette notation par la suite.  Evidemment, l'exposant m
lui-même prête à critique, car  c'est une pratique courante en géométri\unsure{e}
algébrique de désigner par un exp osant la dimension des variétés qui
entrent en jeu. Mais co mme nous n'avons jamais fait usage de ce genre
de convention, je pense que nous gardons les mains libres à ce sujet.

On voit aussitôt, sur la construction précédente de $\underline{X}^{(m)}$, que
l'on a un\inserted{P-}isomorphisme canonique
\[ X^{(m)} = \underset{\unsure{1}}{\underline{\mathrm{Grass}}}_{m}\xout{\$}(\overset{\hand{$\vee$}}{\underline{F}}) , \]
où  $F_{X}$\xout{xxxxxx}  ( $\simeq f^{*}(\underline{\Omega}_{P/S})(1)$)   est le noyau de  $\underline{E}_{X} \longrightarrow \mathcal{O}_{X}(1)$\hand{,}

\mspage{67}
